\documentclass{article}

\usepackage{neurips_2026}

\usepackage[utf8]{inputenc}
\usepackage[T1]{fontenc}
\usepackage[numbers,sort&compress]{natbib}
\usepackage{xurl}
\usepackage{hyperref}
\usepackage{booktabs}
\usepackage{amsfonts}
\usepackage{amsmath}
\usepackage{nicefrac}
\usepackage{microtype}
\usepackage{xcolor}
\usepackage{placeins}
\usepackage{graphicx}

\title{Counterfactual Attribution-Guided Repair for Process-Faithful Logical Reasoning in Large Language Models}

\author{Anonymous Authors}

\begin{document}

\maketitle

\begin{abstract}
Large language models are increasingly used for logical reasoning, but evaluation usually asks only whether the final label is correct. This misses whether the model relies on premises that actually support its conclusion. We introduce CAGE, a counterfactual attribution-guided framework that evaluates and repairs logical reasoning by testing proof-breaking and proof-preserving interventions. CAGE is applicable to direct prompting, Logic-LM-style symbolic translation, and Symbolic Chain-of-Thought, and supports conservative and gated repair policies. Across ProofWriter, PrOntoQA, and FOLIO, CAGE-Conservative preserves initially valid labels, repairs invalid outputs, and often improves counterfactual and support metrics for structured base methods. CAGE-Gated yields substantial gains on FOLIO, improving Direct from 44.1\% to 67.2\% and SymbCoT from 47.1\% to 64.2\%, while revealing when repair is unreliable. CAGE provides a unified evaluation of answer correctness and process faithfulness.
\end{abstract}

\section{Introduction}

Logical reasoning is often evaluated as an answer prediction problem: given a set of facts, rules, and a query, a model is scored by whether it predicts \texttt{true}, \texttt{false}, or \texttt{unknown}. This evaluation is necessary but incomplete. In logical reasoning, an answer can be correct for the wrong reason: a model may rely on lexical overlap, a spurious premise, or an unstable heuristic while still matching the gold label. Conversely, a model may produce a plausible natural-language explanation that does not correspond to the premises that actually justify the answer. Therefore, answer accuracy alone cannot determine whether an LLM performs process-faithful reasoning.

Recent work has substantially improved LLM logical reasoning by strengthening answer-generation pipelines. Logic-LM translates natural-language problems into symbolic programs and uses external solvers to improve final predictions \cite{pan2023logiclm}. Symbolic Chain-of-Thought (SymbCoT) elicits symbolic intermediate reasoning steps to make the reasoning process more structured \cite{xu2025faithful}. VeriCoT-style validation checks completed reasoning traces for logical consistency \cite{vericot2025}. These methods are important advances, but their primary success criterion remains whether the final answer is correct. They do not fully answer a complementary question: \emph{does the model's stated reasoning depend on the premises that causally justify its answer?}

We argue that logical-reasoning evaluation should include \textbf{reason accuracy}, or process faithfulness, in addition to answer accuracy. A reason is faithful only if it behaves causally: changing a premise that is necessary for the proof should change or invalidate the answer, while changing an irrelevant distractor should not. This criterion naturally suggests a counterfactual evaluation protocol. Rather than only asking whether a model solved the original problem, we also ask whether its prediction remains stable under proof-preserving interventions and changes appropriately under proof-breaking or contradiction interventions.

To operationalize this idea, we propose \textbf{Counterfactual Attribution-Guided Evaluation and Repair (CAGE)}. CAGE wraps an existing reasoning method and probes its prediction with counterfactual variants of the original example. For ProofWriter and PrOntoQA-style data, CAGE constructs proof-breaking, proof-preserving, support-shift, contradiction, and distractor interventions. For FOLIO, where the examples do not expose the same structured counterfactual format, CAGE uses transfer probes over natural-language premises, formal representations, negated conclusions, and verifier views. These probes produce diagnostic signals about whether a base method's answer and claimed support are causally aligned with the problem.

CAGE is designed to be method-agnostic. We apply it to three base methods: direct prompting, Logic-LM-style symbolic translation, and SymbCoT. The wrapper first extracts the base method's label and claimed support, then runs counterfactual diagnostics, and finally decides whether repair is safe. A key empirical lesson is that counterfactual feedback should not always be used to rewrite answers. Early pilot experiments showed that aggressive free-form repair can reduce answer accuracy. We therefore introduce a conservative repair policy: if the base method returns a valid label, CAGE preserves the label and records diagnostics; repair is triggered only when the initial output is invalid. This policy targets reason accuracy while limiting unnecessary damage to answer accuracy.

Our contributions are as follows. First, we propose a counterfactual evaluation standard for reason accuracy in LLM logical reasoning. Second, we introduce CAGE, a method-agnostic framework that combines counterfactual attribution diagnostics with conservative and gated repair. Third, we instantiate CAGE for direct prompting, Logic-LM-style reasoning, and SymbCoT across ProofWriter, PrOntoQA, and FOLIO. Fourth, we show that, in our experiments, conservative CAGE eliminates invalid outputs while preserving initially valid labels, and that gated CAGE can produce large gains when repairs are reliable, especially on FOLIO.

\begin{figure}[t]
\centering
\includegraphics[width=0.9\linewidth]{cage_overview.png}
\caption{Overview of CAGE. A base reasoning method first predicts an answer label and claimed support; CAGE then applies counterfactual attribution diagnostics and chooses between conservative preservation and gated repair. The framework reports both final-answer quality and process-faithfulness signals, including invalid-output rate, attribution quality, counterfactual robustness, and repair reliability.}
\label{fig:cage-overview}
\end{figure}

\section{Related work}

\paragraph{Logical reasoning and neuro-symbolic methods}
ProofWriter, PrOntoQA, and FOLIO evaluate natural-language logical reasoning with labels such as \texttt{true}, \texttt{false}, and \texttt{unknown} \cite{clark2020transformers,saparov2022proofwriter,saparov2022prontoqa,hwang2023folio}. They are useful for our setting because their inputs contain explicit premises rather than only free-form questions. Chain-of-thought prompting improves answer generation by eliciting intermediate reasoning, but such reasoning can be persuasive without being causally faithful \cite{wei2022chain,wang2023self}. Neuro-symbolic systems add further structure: Logic-LM translates language into symbolic programs and uses solvers, while SymbCoT elicits symbolic intermediate reasoning \cite{pan2023logiclm,xu2025faithful}. CAGE is complementary to these methods because it does not replace the base reasoner; it wraps the base prediction and tests whether the prediction is causally supported by the premises.
\paragraph{Faithfulness and counterfactual evaluation}
Faithfulness work argues that plausible explanations are not enough: a faithful reason should change when necessary evidence changes and remain stable under irrelevant perturbations \cite{jacovi2020towards,wiegreffe2019attention,jain2019attention,chen2023faithful}. Counterfactual and causal evaluation provide a natural way to test this dependence by intervening on inputs or rationales and measuring whether predictions change in the expected direction \cite{kaushik2020learning,wu2023counterfactual}. Existing work often evaluates explanations or rationales after the fact; CAGE instead makes premise-level interventions part of the logical-reasoning pipeline and uses their outcomes to decide whether repair is allowed.
\paragraph{Self-repair and verifier-guided reasoning}
A related line of work uses verifiers, critics, or self-refinement loops to improve model outputs after an initial answer \cite{madaan2023self,shinn2023reflexion,vericot2025}. These methods show that feedback can improve reasoning, but unconstrained repair can also introduce new errors. CAGE differs by separating diagnostic evidence from answer rewriting: conservative CAGE records diagnostics without changing valid labels, while gated CAGE accepts a changed answer only when diagnostic failures provide a strict alternative-label majority. We therefore evaluate repair quality explicitly, including accepted precision and harmful accepted repairs.

\section{Task formulation}

Each example consists of a context $x$, a query $q$, and a gold label $y \in \{\texttt{true}, \texttt{false}, \texttt{unknown}\}$. A method returns a prediction
\[
  \hat{y}, S, r = f(x, q),
\]
where $\hat{y}$ is the answer label, $S$ is a set of claimed causal premises, and $r$ is the raw model response or trace. Standard accuracy evaluates whether $\hat{y}=y$. We additionally evaluate whether $S$ behaves like a causal support set.

For a counterfactual example $x'$ derived from $x$, let $\rho(x,x')$ denote the expected relation between the original and counterfactual answers. If a support premise is removed or contradicted, $\rho$ often expects the answer to change or become unknown. If an irrelevant distractor is changed, $\rho$ expects the answer to be preserved. A counterfactual diagnostic passes when the model's prediction on $x'$ matches the expected relation. We report these diagnostics as a verifier counterfactual acceptance rate (VCAR) and as failure-type breakdowns.

\section{Method}

\subsection{Base reasoning methods}
CAGE wraps three base reasoning families. \textbf{Direct prompting} receives the original context and query and returns a natural-language response from which we normalize the final answer label. It does not expose a formal proof trace, so claimed support is extracted, when present, from cited premise identifiers or premise mentions in the response. \textbf{Logic-LM-style reasoning} follows a neuro-symbolic interface: the model translates the natural-language problem into a symbolic representation, an external solver or deterministic checker derives a label when the translation is executable, and support is taken from the premises used by the symbolic derivation when available. Outputs that cannot be parsed into a valid label are marked \texttt{invalid}. \textbf{SymbCoT} elicits symbolic intermediate reasoning steps before the final answer. We parse both the final label and the symbolic trace, then map referenced steps or premises back to the original support set when the trace exposes such links. Across all three methods, CAGE receives the same normalized tuple $(\hat{y},S,r)$: a label, a claimed support set, and the raw response or trace.

FOLIO requires a different diagnostic interface because the examples do not expose the same structured proof counterfactuals. For FOLIO, CAGE uses transfer probes over four views of the problem: natural-language premises only, formal representation only, a negated-conclusion probe, and an answer-verifier probe. These probes are treated as consistency diagnostics rather than gold support annotations. A disagreement or conflict among these views is evidence that the initial answer may be unstable, but not by itself sufficient to force a repair.
\subsection{Repair policies}
CAGE separates diagnostic collection from answer rewriting. This distinction is important because counterfactual feedback can reveal process failures even when changing the final label would hurt accuracy. We therefore instantiate two repair policies.
\paragraph{CAGE-Conservative}
The conservative policy is the main safe wrapper. It runs the base method, extracts the initial label and claimed premises, runs diagnostics, and stores the full diagnostic trace. If the initial label is valid, i.e., $\hat{y}\in\{\texttt{true},\texttt{false},\texttt{unknown}\}$, the wrapper preserves the label exactly. Repair is triggered only when the initial label is invalid. In that case, CAGE prompts the model with the original problem, the initial raw response, the claimed premises, and the failed diagnostics, and accepts the repaired answer only if it is a valid label. Thus any accuracy gain under conservative CAGE comes only from converting invalid outputs into valid labels; the policy does not correct wrong-but-valid predictions.
\paragraph{CAGE-Gated}
The gated policy is an accuracy-seeking variant. It may repair a valid initial label, but only when failed diagnostics provide a strict alternative-label majority. For structured datasets, the gate collects failed diagnostics whose failure type is not \texttt{ok}; for FOLIO, it collects probes marked as disagreement or conflict. A repair target is formed only if at least two diagnostics fail and more than half of their valid alternative labels agree on the same label different from the initial label. The repair prompt is then run, but the repaired label is accepted only if it exactly matches the gated target. Otherwise, the wrapper keeps the original prediction. This acceptance rule is designed to prevent arbitrary free-form repair while still allowing diagnostic evidence to revise unstable answers.
Both policies write the same kind of audit trace: the base method name, raw base output, initial answer, claimed premises, diagnostic records, repair trigger reason, optional gated target, repair acceptance status, and final answer. These traces support the repair-quality metrics reported in Section~\ref{sec:repair-quality} and Table~\ref{tab:repair-quality}: trigger rate, accepted changes, accepted precision, and harmful accepted repairs.


\section{Experiments}

\subsection{Datasets, variants, and metrics}
\label{sec:datasets-metrics}

We evaluate on ProofWriter, PrOntoQA, and FOLIO. The structured runs use 500 ProofWriter examples and 499 parseable PrOntoQA examples; FOLIO uses the full 204-example validation set. We compare three base reasoning families---direct prompting, Logic-LM-style symbolic translation, and SymbCoT---with two CAGE variants. \textbf{CAGE-Conservative} runs diagnostics but preserves any valid base label, repairing only invalid labels. \textbf{CAGE-Gated} is an accuracy-seeking variant that may change a valid label only when failed diagnostics provide a strict alternative-label majority and the repair agrees with that target.

We report answer accuracy, invalid-label rate, attribution F1, exact support match, counterfactual accuracy, counterfactual consistency, and repair quality. Attribution metrics are available for ProofWriter and PrOntoQA, whose examples expose support annotations. FOLIO does not expose the same support fields in this experiment layout, so for FOLIO we emphasize answer accuracy, invalid rate, and repair-quality statistics.

\subsection{Implementation details}

All LLM calls use a provider-agnostic interface, \texttt{call\_llm(prompt, max\_tokens, json\_output, schema)}, rather than provider-specific code inside CAGE. The reported full-scale runs use DeepSeek through this interface. Unless overridden by \texttt{DEEPSEEK\_MODEL} or \texttt{CF\_REASONING\_MODEL}, the model is \texttt{deepseek-chat}; decoding uses temperature 0, one prompt attempt per example, and the configured maximum output length \texttt{CF\_REASONING\_MAX\_TOKENS} (512 by default). Prompts, parsing rules, the conservative policy, and the gated-repair threshold are fixed before evaluation: gated repair requires at least two failed diagnostics and a strict majority for one alternative label. We do not tune thresholds on the FOLIO validation labels; the FOLIO transfer probes are used as diagnostic signals and are not equivalent to the structured proof counterfactuals for ProofWriter and PrOntoQA. Labels are normalized to \texttt{true}, \texttt{false}, \texttt{unknown}, or \texttt{invalid}. For structured counterfactual evaluation, we generate two counterfactual examples per original example. Thus the ProofWriter and PrOntoQA counterfactual result files contain 1000 and 998 counterfactual rows, respectively.

\section{Results}

\subsection{Answer accuracy and invalid-output repair}

Table~\ref{tab:main-results} summarizes final answer accuracy and invalid-label rates. Conservative CAGE is the safest main wrapper: it removes invalid outputs while preserving valid base labels. Therefore, the conservative gains in this table should be read as invalid-output repairs, not as corrections to initially valid but wrong labels. On ProofWriter, direct prompting has an inferred base accuracy of 45.6\% and an invalid rate of 7.2\%; Direct+CAGE raises accuracy to 49.2\% and reduces invalid outputs to zero. On PrOntoQA, Direct+CAGE similarly raises inferred direct accuracy from 44.9\% to 46.1\% while eliminating a 2.4\% invalid rate.

\begin{table}[t]
\caption{Answer accuracy and invalid-label rate by dataset. Each subtable reports accuracy (\%) and invalid-label rate (\%) for one dataset. ProofWriter and PrOntoQA use 500 and 499 examples, respectively; FOLIO uses the full 204-example validation set. CAGE denotes the conservative variant; Gated denotes the diagnostic-gated repair variant.}
\label{tab:main-results}
\centering
\vspace{0.5em}
\scriptsize
\setlength{\tabcolsep}{2.5pt}

\begin{minipage}[t]{0.30\linewidth}
\centering
\textbf{FOLIO-204}\\[0.2em]
\begin{tabular}{lrr}
\toprule
Method & Acc. (\%) & Invalid (\%) \\
\midrule
Direct & 44.1 & 0.0 \\
Direct+CAGE & 44.1 & 0.0 \\
Direct+Gated & \textbf{67.2} & 0.0 \\
\midrule
LogicLM & 63.7 & 0.0 \\
LogicLM+CAGE & 63.7 & 0.0 \\
LogicLM+Gated & \textbf{67.6} & 0.0 \\
\midrule
SymbCoT & 47.1 & 1.0 \\
SymbCoT+CAGE & 48.0 & 0.0 \\
SymbCoT+Gated & \textbf{64.2} & 0.0 \\
\bottomrule
\end{tabular}
\end{minipage}
\hspace{0.015\linewidth}
\begin{minipage}[t]{0.30\linewidth}
\centering
\textbf{PrOntoQA-499}\\[0.2em]
\begin{tabular}{lrr}
\toprule
Method & Acc. (\%) & Invalid (\%) \\
\midrule
Direct & 44.9 & 2.4 \\
Direct+CAGE & 46.1 & 0.0 \\
Direct+Gated & 46.5 & 0.0 \\
\midrule
LogicLM & \textbf{62.7} & 0.0 \\
LogicLM+CAGE & \textbf{62.7} & 0.0 \\
LogicLM+Gated & 60.5 & 0.0 \\
\midrule
SymbCoT & 49.1 & 0.0 \\
SymbCoT+CAGE & 49.1 & 0.0 \\
SymbCoT+Gated & \textbf{52.7} & 0.0 \\
\bottomrule
\end{tabular}
\end{minipage}
\hspace{0.015\linewidth}
\begin{minipage}[t]{0.30\linewidth}
\centering
\textbf{ProofWriter-500}\\[0.2em]
\begin{tabular}{lrr}
\toprule
Method & Acc. (\%) & Invalid (\%) \\
\midrule
Direct & 45.6 & 7.2 \\
Direct+CAGE & 49.2 & 0.0 \\
Direct+Gated & 49.2 & 0.0 \\
\midrule
LogicLM & 77.0 & 0.6 \\
LogicLM+CAGE & \textbf{77.2} & 0.0 \\
LogicLM+Gated & 74.4 & 0.0 \\
\midrule
SymbCoT & 50.8 & 0.0 \\
SymbCoT+CAGE & 50.8 & 0.0 \\
SymbCoT+Gated & \textbf{52.8} & 0.0 \\
\bottomrule
\end{tabular}
\end{minipage}

\vspace{-0.5em}
\end{table}

\begin{figure}[t]
\centering
\includegraphics[width=0.85\linewidth]{accuracy_ranked_by_dataset.png}
\caption{Accuracy ranking by dataset. Each panel ranks direct prompting, LogicLM, SymbCoT, and their CAGE variants on one dataset. Gated CAGE gives the largest gains on FOLIO, while LogicLM remains the strongest base method on ProofWriter and PrOntoQA.}
\label{fig:accuracy-ranked}
\end{figure}

The gated variant is not uniformly safe, but it reveals that diagnostic agreement can produce large gains in suitable settings. It is strongest on FOLIO: Direct+Gated improves from 44.1\% to 67.2\%, SymbCoT+Gated from 47.1\% to 64.2\%, and LogicLM+Gated from 63.7\% to 67.6\%. These large changes should be interpreted cautiously: FOLIO uses transfer probes over alternative views rather than the structured proof-preserving and proof-breaking counterfactuals used for the other datasets. The repair audit shows that Direct+Gated accepts 65 changes, with 55 improvements and 8 harms, while SymbCoT+Gated accepts 47 changes, with 39 improvements and 4 harms (Table~\ref{tab:repair-quality}). Thus the gains reflect the behavior of a fixed diagnostic gate on the full validation set, not a theoretical guarantee or a claim that FOLIO probes are equivalent to gold proof interventions. Gated CAGE also improves SymbCoT on PrOntoQA from 49.1\% to 52.7\% and on ProofWriter from 50.8\% to 52.8\%. However, it harms LogicLM on the two structured datasets, reducing ProofWriter accuracy from 77.0\% to 74.4\% and PrOntoQA accuracy from 62.7\% to 60.5\%. We therefore treat CAGE-Gated as an enhanced ablation rather than the main method.

\subsection{Support and counterfactual faithfulness}

Table~\ref{tab:faithfulness} reports support and counterfactual metrics for the structured datasets. Conservative CAGE provides small but consistent improvements for direct prompting in counterfactual robustness: ProofWriter Direct+CAGE reaches 62.2\% counterfactual accuracy, 67.2\% counterfactual consistency, and 81.8\% label-change accuracy; PrOntoQA Direct+CAGE reaches 59.5\%, 65.5\%, and 74.2\%, respectively. LogicLM remains the strongest method on counterfactual metrics, especially on PrOntoQA, where it achieves 80.4\% counterfactual accuracy and 83.0\% consistency.

\begin{table}[t]
\caption{Faithfulness metrics for structured datasets. Attr-F1 (\%) and exact support (\%) measure support recovery on original examples. CF Acc. (\%), CF Cons. (\%), and LC Acc. (\%) denote counterfactual accuracy, counterfactual consistency, and label-change accuracy on generated counterfactual examples. Base Direct rows are omitted because direct prompting does not produce a separate normalized support trace before CAGE extraction; Direct+CAGE is the first row with extracted support. LogicLM+Gated rows are omitted because gated repair changes labels but does not add a new symbolic support extraction pass beyond the LogicLM/CAGE trace. SymbCoT+CAGE rows match the SymbCoT trace when the conservative wrapper preserves valid labels and are therefore omitted to save space.}
\label{tab:faithfulness}
\centering
\vspace{0.5em}
\scriptsize
\begin{tabular}{llrrrrr}
\toprule
Dataset & Method & Attr-F1 (\%) & Exact (\%) & CF Acc. (\%) & CF Cons. (\%) & LC Acc. (\%) \\
\midrule
ProofWriter & Direct+CAGE & 5.7 & 4.0 & 62.2 & 67.2 & 81.8 \\
ProofWriter & Direct+Gated & 9.5 & 6.8 & 62.2 & 67.2 & 81.8 \\
ProofWriter & LogicLM & 65.7 & 52.4 & 72.4 & 77.4 & 79.4 \\
ProofWriter & LogicLM+CAGE & 65.8 & 52.4 & 72.6 & 77.5 & 79.6 \\
ProofWriter & SymbCoT & 0.0 & 0.0 & 43.8 & 57.4 & 31.2 \\
ProofWriter & SymbCoT+Gated & 6.4 & 4.0 & 43.8 & 57.4 & 31.2 \\
\midrule
PrOntoQA & Direct+CAGE & 2.0 & 0.4 & 59.5 & 65.5 & 74.2 \\
PrOntoQA & Direct+Gated & 2.7 & 0.8 & 59.5 & 65.5 & 74.2 \\
PrOntoQA & LogicLM & 67.8 & 36.1 & 80.4 & 83.0 & 89.6 \\
PrOntoQA & LogicLM+CAGE & 67.8 & 36.1 & 80.4 & 83.0 & 89.6 \\
PrOntoQA & SymbCoT & 0.0 & 0.0 & 26.1 & 51.2 & 1.2 \\
PrOntoQA & SymbCoT+Gated & 10.0 & 1.6 & 26.1 & 51.2 & 1.2 \\
\bottomrule
\end{tabular}
\end{table}

\subsection{Repair-quality analysis}
\label{sec:repair-quality}

Table~\ref{tab:repair-quality} separates diagnostic triggering from accepted repairs. This analysis is crucial because gated repair is high-reward but higher-risk. On FOLIO, accepted repairs are highly reliable: Direct+Gated accepts 65 changes, of which 55 improve accuracy and 8 hurt it; SymbCoT+Gated accepts 47 changes, with 39 improvements and 4 harms. In contrast, LogicLM+Gated performs poorly on the structured datasets: on ProofWriter it accepts 20 repairs but only 3 improve accuracy, and on PrOntoQA it accepts 13 repairs but only 1 improves accuracy.

\begin{table}[t]
\caption{Repair-quality audit for gated CAGE. Accepted precision is the percentage of accepted changes that improve accuracy. Unknown-det. wrong counts cases where an initial \texttt{unknown} was changed to \texttt{true} or \texttt{false} incorrectly.}
\label{tab:repair-quality}
\centering
\vspace{0.5em}
\scriptsize
\begin{tabular}{llrrrrr}
\toprule
Dataset & Method & Accepted & Improved & Hurt & Precision (\%) & Unknown-det. wrong \\
\midrule
ProofWriter & Direct+Gated & 61 & 30 & 12 & 49.2 & 0 \\
ProofWriter & LogicLM+Gated & 20 & 3 & 16 & 15.0 & 0 \\
ProofWriter & SymbCoT+Gated & 44 & 27 & 17 & 61.4 & 0 \\
PrOntoQA & Direct+Gated & 16 & 9 & 1 & 56.2 & 0 \\
PrOntoQA & LogicLM+Gated & 13 & 1 & 12 & 7.7 & 0 \\
PrOntoQA & SymbCoT+Gated & 66 & 41 & 23 & 62.1 & 0 \\
FOLIO & Direct+Gated & 65 & 55 & 8 & \textbf{84.6} & 8 \\
FOLIO & LogicLM+Gated & 14 & 11 & 3 & 78.6 & 3 \\
FOLIO & SymbCoT+Gated & 47 & 39 & 4 & \textbf{83.0} & 4 \\
\bottomrule
\end{tabular}
\end{table}

\begin{figure}[t]
\centering
\includegraphics[width=0.75\linewidth]{repair_gain_harm_balance.png}
\caption{Gain--harm balance of gated repair. Green bars show the share of examples changed from wrong to correct, while red bars show correct predictions changed to wrong. The figure illustrates why CAGE-Gated is useful as an enhanced repair variant but should not replace the conservative wrapper as the main safe method.}
\label{fig:repair-gain-harm}
\end{figure}

\FloatBarrier
\section{Discussion}

The results support a two-variant view of CAGE. Conservative CAGE is best understood as a safe process-faithfulness layer. It does not claim to improve every valid answer; instead, it records causal diagnostics and repairs invalid outputs without changing valid base labels. This behavior is useful for settings where answer preservation matters. Gated CAGE, by contrast, turns diagnostic agreement into answer repair. It can be highly effective, as shown by the large FOLIO gains, but it can also over-repair strong baselines. The contrast between FOLIO and LogicLM on structured datasets shows that repair decisions should depend on both diagnostic strength and base-method reliability.

The counterfactual results also clarify what answer accuracy hides. LogicLM is strongest on support and counterfactual metrics, while SymbCoT underperforms on counterfactual proof-breaking despite its structured traces. This suggests that symbolic-looking traces are not automatically causally faithful. CAGE makes this distinction measurable by testing whether predictions behave correctly under proof-relevant interventions.

\section{Limitations}

The main limitation is repair stability. Conservative CAGE is intentionally safe because it preserves valid base labels, but this also means it cannot fix wrong-but-valid predictions. Gated CAGE can improve weaker or less stable methods, especially on FOLIO and SymbCoT, but its benefits are not uniform across base reasoners. In particular, LogicLM already produces comparatively strong and stable predictions, and gated repair can overrule them incorrectly, reducing accuracy on ProofWriter and PrOntoQA. This suggests that diagnostic agreement alone is not always sufficient for safe answer revision; future variants should calibrate repair thresholds to the reliability of the underlying base method. FOLIO also uses transfer probes rather than the structured proof counterfactuals available for ProofWriter and PrOntoQA, so its diagnostic signals are not directly identical across datasets.

\section{Conclusion}

We introduced CAGE, a method-agnostic framework for counterfactual attribution-guided diagnosis and repair in LLM logical reasoning. Conservative CAGE provides a safe wrapper that preserves initially valid answers, repairs invalid labels, and records process-faithfulness diagnostics. Gated CAGE shows that the same diagnostics can produce large accuracy gains when repair is reliable, especially on FOLIO and on weaker base methods, but it is not universally safe. Together, these results argue that logical-reasoning evaluation should report not only final answer accuracy, but also invalid-output rates, support attribution, counterfactual robustness, and repair quality.

\FloatBarrier
\bibliographystyle{unsrtnat}
\bibliography{references}

\appendix

\section{Additional result visualizations}
\begin{figure}[!htbp]
\centering
\includegraphics[width=0.8\linewidth]{accuracy_grouped_overview.png}
\caption{Complete grouped accuracy comparison across datasets and method families. For each dataset, the figure compares the base method, CAGE-Conservative, and CAGE-Gated variants for direct prompting, LogicLM, and SymbCoT.}
\label{fig:accuracy-grouped}
\end{figure}

\FloatBarrier
\section{Reproducibility notes}

The smoke experiments are only local sanity checks before full-scale API runs; no smoke-run numbers are reported as main results. The reported main results use the full evaluation described in Section~\ref{sec:datasets-metrics}: 500 ProofWriter examples, 499 parseable PrOntoQA examples, and the 204-example FOLIO validation set.
\begin{itemize}
\item Provider and model: runs use DeepSeek through the provider-agnostic LLM interface with \path|CF_REASONING_LLM_PROVIDER=deepseek|, \path|DEEPSEEK_API_KEY|, and optionally \path|DEEPSEEK_MODEL|; by default, the implementation uses \path|deepseek-chat|.
\item Decoding and runs: decoding uses temperature 0, a single prompt attempt per method per example, deterministic label normalization, and \path|CF_REASONING_MAX_TOKENS=512| unless otherwise configured.
\item Evaluation protocol: all method prompts and parsing rules are fixed before running the reported evaluations. Conservative CAGE preserves initially valid labels by construction, and gated CAGE uses the same pre-specified threshold in all datasets: at least two failed diagnostics and a strict majority supporting one alternative label. FOLIO validation labels are used only for evaluation, not to tune the gate threshold.
\end{itemize}

\end{document}

